\section{From twisted differentials to effective divisors}
\label{sec:differentials}
Theorem~\ref{thm:conversion} accepts any fixed pseudoeffective error.
We now construct one from differential forms, so that the invariant
forms obtained in Section~\ref{sec:index} can be used after descent.
The construction has two steps: take a fixed determinant line of a generic
span, then control the sign of that line by generic nefness.

\begin{theorem}[Cotangent subsheaves]\label{thm:subsheaf}
Let \(S\) be smooth projective with \(-K_S\) nef.
If a line bundle \(M\) is a saturated subsheaf of
\((\Omega_S^1)^{\otimes b}\), where \(b>0\), then \(-M\) is
pseudoeffective.
\end{theorem}

This is the smooth, zero-boundary, rank-one case of
\cite[Theorem~4.1]{LMPTX2023}; Ou's generic-nefness theorem
\cite[Theorem~1.4]{Ou2023} is an antecedent of the cotangent result.

\begin{lemma}[Determinants of twisted differentials]\label{lem:determinant}
Let \(S\) be smooth connected projective with \(-K_S=L\) nef.
Suppose that, for a fixed \(p>0\),
\[
H^0(S,\Omega_S^p\otimes\OO_S(mL))\ne0
\]
for arbitrarily large positive integers \(m\).
Then there is a pseudoeffective Cartier divisor \(D\) and a strictly
increasing sequence \(d_j>0\) with effective integral divisors
\[
N_j\sim d_jL-D.
\]
\end{lemma}

The generic-span construction below adapts
\cite[Lemma~4.1]{LP2018} and \cite[Lemma~5.1]{LMPTX2023}. We give the
proof for the present unbounded positive sequence, including the
numerically trivial case excluded by the numerical-nontriviality
premises of those two cited lemmas, and include the saturation argument needed
for the negative first-Chern-class conclusion of
\cite[Theorem~4.1]{LMPTX2023}.

\begin{proof}
Choose one nonzero section at each of an unbounded set of exponents.
Over \(\C(S)\), trivialize \(L\) rationally and let \(W\) be the
span of the chosen vectors. Take a basis \(u_1,\ldots,u_r\) of \(W\)
from among them. For infinitely many unbounded exponents, the corresponding
vector has nonzero coefficient along one fixed \(u_a\). Wedging that
section with the other \(r-1\) basis sections gives a nonzero section
of
\[
\bigwedge\nolimits^r\Omega_S^p\otimes\OO_S(d_jL),
\]
where, if the fixed basis sections have exponents \(a_1,\ldots,a_r\),
the new exponent is \(d_j=m_j+\sum_{i\ne a}a_i\).
The sum is fixed, so the \(d_j\) are unbounded and may be taken strictly increasing.
Every such wedge spans the same line \(\bigwedge^r W\).

Let \(M\) be its saturation in \(\bigwedge^r\Omega_S^p\).
A saturated rank-one subsheaf of a vector bundle on a smooth variety
is reflexive: its double dual maps into that vector bundle by extension
across codimension two, and any enlargement would give torsion in the
quotient. Hence \(M\) is a line bundle. Each wedge factors through
\(M(d_jL)\), since its image in the torsion-free quotient is generically
zero.

In characteristic zero, exterior powers are direct summands of tensor
powers. Thus \(\bigwedge^r\Omega_S^p\) is a direct summand of
\((\Omega_S^1)^{\otimes rp}\), and \(M\) remains saturated in that
tensor power: the additional summand is a vector bundle, so the enlarged
quotient is still torsion-free. Theorem~\ref{thm:subsheaf} makes \(D=-M\)
pseudoeffective. The zero divisors of the induced nonzero sections of
\(M(d_jL)\) are the required \(N_j\).
\end{proof}

\begin{corollary}[Conversion of twisted differentials]\label{cor:conversion}
Let \(S\) be smooth connected projective with smoothly semipositive
\(L=-K_S\). If, for one fixed \(p\ge0\),
\[
H^0(S,\Omega_S^p\otimes\OO_S(mL))\ne0
\]
for arbitrarily large positive integers \(m\), then some positive
multiple of \(L\) has a nonzero section.
\end{corollary}

\begin{proof}
For \(p=0\) this is immediate. For \(p>0\), apply
Lemma~\ref{lem:determinant} and Theorem~\ref{thm:conversion}.
\end{proof}
